\section{Extensions and Alternatives}

\subsection{Non-Abelian Gauge Theories}

Replacing $\mathbb{Z}_8$ with a non-abelian group (e.g., SU(2)$_k$, D$_8$) introduces non-abelian anyons. A single non-abelian defect can encode a superposition of multiple abelian charges.

\begin{itemize}
    \item \textbf{Fibonacci anyons} (SU(2)$_3$): One anyon can replace an entire cluster of $\mathbb{Z}_8$ defects.
    \item \textbf{Ising anyons} (Z$_2$ gauge theory with fermions): Can encode T gates more efficiently.
\end{itemize}

\textit{Hypothesis:} In non-abelian theories, defect complexity may be polynomial for BQP states.

\subsection{Continuous Connections on Spin-Foams}

Instead of discrete fluxes, we can use a continuous connection $A_\mu(x)$ on a 4D spin-foam. The state is specified by holonomies along a finite set of curves.

For integrable models (Ashtekar-Lewandowski), holonomies completely determine the wavefunction. Expectation values become path integrals that may be polynomial if the number of holonomy variables is bounded.

\subsection{Hybrid Constructions}

We can combine abelian and non-abelian elements:
\begin{itemize}
    \item Use $\mathbb{Z}_8$ defects for Clifford parts of circuits.
    \item Use non-abelian defects for T gates and magic states.
\end{itemize}

This may yield polynomial description and evaluation for a broader class of states.

\subsection{Open Questions}

\begin{enumerate}
    \item Does there exist a non-abelian gauge theory with polynomial defect complexity for all BQP states?
    \item Can continuous connections be discretized with polynomial error?
    \item What is the precise relationship between defect complexity and measures of magic (stabilizer rank, mana, robustness)?
\end{enumerate}
